% ============================================================
%  Chapter 3: Method
%  File: chapters/03_method.tex
%  Bibliography: chapters/refs_03_method.bib
% ============================================================

\chapter{Method}
\label{ch:method}

This chapter describes the \modnet{} architecture and the evolutionary
algorithm that discovers its structure. We begin with the design
principles, then describe the node model, the layer-free connectivity,
the recurrence mechanism, the fitness function, and the four
evolutionary operators: mutation, crossover, birth, and pruning.

% ------------------------------------------------------------
\section{Design Principles}
\label{sec:method-principles}
% ------------------------------------------------------------

\modnet{} is built on five design principles that distinguish it from
existing neuroevolutionary architectures.

\begin{enumerate}
  \item \textbf{No predefined layers.} There is no distinction between
    input, hidden, and output layers. Every node is a general-purpose
    computational unit that may receive input from any other node and
    may send output to any other node.

  \item \textbf{No modular grouping.} Unlike our earlier modular variant,
    \modnet{} has no modules at all. The nodes form a single flat set.
    This is a deliberate departure: we want to test whether emergent
    structure can be obtained \emph{without} any modular scaffold.

  \item \textbf{Recurrence as a first-class primitive.} Any node may
    connect to itself or to any other node, regardless of index. The
    network is therefore a general directed graph, not a directed acyclic
    graph. Recurrence is not added as a special case; it is a natural
    consequence of unrestricted connectivity.

  \item \textbf{Growth through addition and pruning.} The network begins
    with a small number of nodes and grows as needed. New nodes are
    appended to the network, and existing nodes are removed when they
    become inactive or isolated. Growth is bidirectional.

  \item \textbf{Gradient-free learning.} No gradient information is used
    at any stage. All learning is performed by a compact evolutionary
    algorithm operating on the topology and weights simultaneously.
\end{enumerate}

These principles are motivated by the observation that biological neural
networks exhibit none of the structural regularities that are imposed on
artificial networks by design. The brain has no layers, no modular
scaffold, and no global error signal; it discovers its structure through
developmental and evolutionary processes
\citep{gershenson2012guiding, plantec2024evolving}. \modnet{} is an
attempt to capture these properties in a minimal computational model.

% ------------------------------------------------------------
\section{Node Model}
\label{sec:method-node}
% ------------------------------------------------------------

A node in \modnet{} is a threshold unit with a weighted sum of inputs.
Let $N$ be the total number of nodes in the network. For a node
$i \in \{0, \ldots, N-1\}$, let $\mathcal{S}_i$ be the set of source
indices from which node $i$ receives input. Each source
$s \in \mathcal{S}_i$ may be either an external input (denoted by a
negative index $-k$ for input $k$) or another node (denoted by its
index $j \geq 0$). Node $i$ is also equipped with a bias $b_i$.

At each time step $t$, the state of node $i$ is updated according to

\begin{equation}
  z_i^{(t)} = b_i + \sum_{s \in \mathcal{S}_i} w_{i,s} \, a_s^{(t)},
  \label{eq:method-pre-activation}
\end{equation}

\begin{equation}
  a_i^{(t+1)} =
  \begin{cases}
    1 & \text{if } z_i^{(t)} \geq 0, \\
    0 & \text{otherwise},
  \end{cases}
  \label{eq:method-activation}
\end{equation}

where $a_s^{(t)}$ is the activation of source $s$ at time $t$.
External inputs are held constant across time steps: for an input
index $-k$, we have $a_{-k}^{(t)} = x_k$ for all $t$.

The activation function is the Heaviside step function, which is
non-differentiable. This is why we use evolutionary rather than
gradient-based learning: a differentiable surrogate is not required.

\begin{figure}[h]
  \centering
  \begin{tikzpicture}[
    node distance=1.2cm,
    every node/.style={font=\small},
    neuron/.style={circle, draw, minimum size=1cm, thick},
    input/.style={rectangle, draw, minimum size=0.7cm, fill=gray!20},
    arrow/.style={-{Latex[length=2mm]}, thick},
  ]
    \node[input] (x1) at (-3, 1.5) {$x_1$};
    \node[input] (x2) at (-3, 0)   {$x_2$};
    \node[input] (x3) at (-3, -1.5) {$x_3$};

    \node[neuron] (n) at (0, 0) {$i$};
    \node[neuron] (out) at (2.5, 0) {$j$};

    \draw[arrow] (x1) -- node[above, sloped, font=\scriptsize] {$w_{i,1}$} (n);
    \draw[arrow] (x2) -- node[above, font=\scriptsize] {$w_{i,2}$} (n);
    \draw[arrow] (x3) -- node[below, sloped, font=\scriptsize] {$w_{i,3}$} (n);
    \draw[arrow] (n) -- (out);

    \draw[arrow] (n) to[out=135, in=45, looseness=8]
      node[above, font=\scriptsize] {$w_{i,i}$} (n);

    \node[font=\scriptsize, below=0.6cm of n]
      {$z_i = b_i + \sum_s w_{i,s} a_s$};
  \end{tikzpicture}
  \caption{Node model. Each node computes a weighted sum of its
    sources plus a bias, then applies a Heaviside step function. The
    self-loop indicates that the node may receive input from itself,
    which introduces recurrence.}
  \label{fig:method-node}
\end{figure}

% ------------------------------------------------------------
\section{Layer-Free Connectivity}
\label{sec:method-connectivity}
% ------------------------------------------------------------

The central architectural choice in \modnet{} is the absence of both
layers and modules. A connection from node $i$ to node $j$ is permitted
if and only if $i \neq j$ is false---that is, self-connections are
explicitly allowed. Formally, the connectivity matrix
$\mathbf{W} \in \mathbb{R}^{N \times N}$ has no structural zeros: any
entry $W_{i,j}$ may be non-zero.

Connections from external inputs are handled separately. We use a
separate input weight matrix $\mathbf{W}_{\text{in}} \in
\mathbb{R}^{N \times K}$, where $K$ is the number of external inputs.
A node $i$ may receive input from any subset of the $K$ inputs, and
this subset is determined by evolution.

\subsection{Output Selection}

The first $n_{\text{out}}$ nodes of the network---that is, nodes with
indices $0, 1, \ldots, n_{\text{out}} - 1$---are designated as
\emph{output nodes}. The output of the network is read from these nodes
after $T$ time steps. There is no separate output layer; the outputs
are simply ordinary nodes whose activations are observed.

During pruning, output nodes are protected: they are never removed,
regardless of their activity or connectivity.

\begin{remark}
The absence of layers has two important consequences. First, the
concept of ``depth'' does not apply: a signal may traverse the network
in cycles of arbitrary length, and information may flow in any
direction. Second, there is no natural topological ordering of nodes;
the index $i$ is purely a labeling convention and does not imply any
temporal or structural priority.
\end{remark}

To ensure that the network has well-defined behavior, we evaluate it for
a fixed number of time steps $T$, which is a hyperparameter. The state
of all nodes is initialized to zero, and after $T$ steps the output is
read from the designated output nodes.

% ------------------------------------------------------------
\section{Recurrence}
\label{sec:method-recurrence}
% ------------------------------------------------------------

Recurrence in \modnet{} arises from two sources: self-connections and
cycles of length greater than one. Both are permitted, and both may
coexist in the same network.

A self-connection on node $i$ is simply a non-zero entry $W_{i,i}$.
Such a connection allows the node to retain information across time
steps: if the node is active at time $t$ and its self-weight is
positive, it may remain active at time $t+1$ even if no other input
changes.

A cycle of length $\ell > 1$ is a sequence of nodes
$i_1 \to i_2 \to \cdots \to i_\ell \to i_1$ where each consecutive
pair has a non-zero connection. Cycles allow information to be
propagated, transformed, and returned to earlier parts of the network.

We quantify recurrence using the number of self-connections:

\begin{equation}
  R = \left| \{i : |W_{i,i}| > \epsilon\} \right|,
  \label{eq:method-recurrence-metric}
\end{equation}

where $\epsilon = 0.05$ is a small threshold below which a weight is
considered inactive. We report $R$ for each evolved network in
Chapter~\ref{ch:results}. As we will show, $R$ is not imposed by the
design; it emerges spontaneously and correlates with the temporal
complexity of the task.

% ------------------------------------------------------------
\section{Fitness Evaluation}
\label{sec:method-fitness}
% ------------------------------------------------------------

Let $\mathcal{D} = \{(\mathbf{x}^{(k)}, \mathbf{y}^{(k)})\}_{k=1}^{K}$
be a set of input--output pairs. For a given network, let
$\hat{\mathbf{y}}^{(k)}$ be the output produced when the network is
initialized to zero and run for $T$ time steps with input
$\mathbf{x}^{(k)}$.

We use two fitness functions depending on the task type, both
formulated so that \emph{higher is always better} and the maximum
attainable value is zero.

\subsection{Boolean tasks}

For Boolean tasks with a single output bit, the fitness is the negative
mean squared error between the predicted scalar and the target scalar:

\begin{equation}
  F_{\text{bool}}(\modnet{}) = -\frac{1}{K} \sum_{k=1}^{K}
    \left( \hat{y}^{(k)} - y^{(k)} \right)^2,
  \label{eq:method-fitness-boolean}
\end{equation}

where $\hat{y}^{(k)}$ is the decoded output (either $0$ or $1$).

For Boolean tasks with multiple output bits, the fitness is the negative
mean squared error computed at the bit level:

\begin{equation}
  F_{\text{bit}}(\modnet{}) = -\frac{1}{K \cdot n_{\text{out}}}
    \sum_{k=1}^{K} \sum_{j=1}^{n_{\text{out}}}
    \left( \hat{b}_j^{(k)} - b_j^{(k)} \right)^2,
  \label{eq:method-fitness-bit}
\end{equation}

where $\hat{b}_j^{(k)}$ and $b_j^{(k)}$ are the $j$-th predicted and
target bits, respectively.

\subsection{Regression tasks}

For regression tasks, the output is a scalar obtained by decoding the
$n_{\text{out}}$ output bits as an integer and normalizing:

\begin{equation}
  \text{decode}(\mathbf{b}) = \frac{1}{2^{n_{\text{out}}} - 1}
    \sum_{j=0}^{n_{\text{out}}-1} b_j \cdot 2^j.
  \label{eq:method-decode}
\end{equation}

The fitness is then the negative mean squared error between the decoded
output and the target value:

\begin{equation}
  F_{\text{reg}}(\modnet{}) = -\frac{1}{K} \sum_{k=1}^{K}
    \left( \text{decode}(\hat{\mathbf{b}}^{(k)}) - y^{(k)} \right)^2.
  \label{eq:method-fitness-regression}
\end{equation}

\begin{remark}
The use of a single scalar fitness for both Boolean and continuous
tasks, with only the decoding step changing, emphasizes that the
architecture does not need to be adapted to the task type. The
selection mechanism is identical in all cases.
\end{remark}

% ------------------------------------------------------------
\section{Evolutionary Algorithm}
\label{sec:method-evolution}
% ------------------------------------------------------------

The evolutionary algorithm in \modnet{} is a compact generational
algorithm with four operators: mutation, crossover, birth, and pruning.
The algorithm maintains a population of $P$ networks. At each
generation, the population is evaluated, sorted by fitness, and used to
produce the next generation.

\begin{algorithm}[h]
  \caption{ModNet Evolution}
  \label{alg:method-evolution}
  \begin{algorithmic}[1]
    \State Initialize population $\mathcal{P} \gets \{N_1, \ldots, N_P\}$
      with random topologies of size $N_0$.
    \State Evaluate $F(N_p)$ for all $p \in \{1, \ldots, P\}$.
    \For{$g = 1, 2, \ldots, G$}
      \State Sort $\mathcal{P}$ by fitness in descending order.
      \State $\mathcal{E} \gets$ top $\lfloor P/5 \rfloor$ networks.
      \State $\mathcal{P}' \gets \mathcal{E}$.
      \While{$|\mathcal{P}'| < P$}
        \State $r \gets \text{Uniform}(0, 1)$.
        \If{$r < p_{\text{cross}}$ and $|\mathcal{E}| \geq 2$}
          \State $N_{\text{new}} \gets \text{Crossover}(N_a, N_b)$.
        \ElsIf{$r < 1 - p_{\text{fresh}}$}
          \State $N_{\text{new}} \gets \text{Mutate}(N_a)$.
        \Else
          \State $N_{\text{new}} \gets$ random new network.
        \EndIf
        \State $\mathcal{P}' \gets \mathcal{P}' \cup \{N_{\text{new}}\}$.
      \EndWhile
      \If{$g \bmod g_{\text{birth}} = 0$}
        \State Apply \textsc{Birth} to top 3 networks.
      \EndIf
      \If{$g \bmod g_{\text{prune}} = 0$}
        \State Apply \textsc{Prune} to top $\lfloor P/5 \rfloor$
          networks.
      \EndIf
      \State $\mathcal{P} \gets \mathcal{P}'$.
    \EndFor
    \State \Return best network in $\mathcal{P}$.
  \end{algorithmic}
\end{algorithm}

\subsection{Mutation}
\label{sec:method-mutation}

The mutation operator modifies the weights, biases, and connectivity of
a network. For each mutation step, a node $i$ is selected uniformly at
random, and one of four operations is applied:

\begin{enumerate}
  \item With probability $0.40$, perturb the weight of a randomly
    selected incoming connection:
    $W_{i,j} \gets W_{i,j} + \mathcal{N}(0, \sigma)$.
  \item With probability $0.20$, perturb the bias:
    $b_i \gets b_i + \mathcal{N}(0, \sigma)$.
  \item With probability $0.20$, perturb the weight of an input
    connection:
    $W^{\text{in}}_{i,k} \gets W^{\text{in}}_{i,k} +
    \mathcal{N}(0, \sigma)$.
  \item With probability $0.20$, either remove an existing connection
    (if the node has more than one active incoming connection) or add a
    new connection with a weight drawn from $\mathcal{N}(0, 1)$.
\end{enumerate}

The mutation scale $\sigma$ is not fixed; it grows when the population
reaches a plateau, according to the schedule

\begin{equation}
  \sigma \gets \min(\sigma_{\max}, \sigma \cdot \gamma)
  \quad \text{when } g_{\text{plateau}} > g_{\text{thresh}},
  \label{eq:method-sigma-adaptation}
\end{equation}

where $g_{\text{plateau}}$ counts generations since the last fitness
improvement, and $g_{\text{thresh}}$ is a hyperparameter. This allows
the population to escape local optima by increasing the exploration
radius when needed.

\subsection{Crossover}
\label{sec:method-crossover}

The crossover operator combines two parent networks into a single
child. It requires that the two parents have the same number of nodes;
if this is not the case, the operator returns \texttt{None} and the
caller falls back to mutation.

For each node $i$, with probability $0.5$, the entire row of
$\mathbf{W}$, the entire row of $\mathbf{W}^{\text{in}}$, and the bias
$b_i$ are copied from the second parent; otherwise they are kept from
the first parent.

This row-wise crossover is a lightweight alternative to the historical
markings used in NEAT \citep{stanley2002evolving}, which would require
tracking the lineage of each connection. Because \modnet{} does not use
historical markings, we restrict crossover to parents of identical
shape. This is a deliberate simplification.

\subsection{Birth}
\label{sec:method-birth}

The birth operator adds a single new node to the network. The new node
is appended at the end of the node list; that is, its index is $N$, the
current total number of nodes.

The new node receives a small number of random incoming connections,
with weights drawn from $\mathcal{N}(0, 1)$, and a random bias. Because
the new node is hidden (not an output), it may connect directly to
external inputs.

Growth is applied periodically (every $g_{\text{birth}}$ generations)
to the top three networks in the population. This ensures that the
network size is allowed to increase when the task demands it, but only
in response to sustained selection pressure.

\subsection{Pruning}
\label{sec:method-prune}

The pruning operator removes nodes that are no longer contributing to
the network's behavior. A node is considered for removal if any of the
following conditions hold:

\begin{enumerate}
  \item \textbf{Dead:} its mean activation across all training samples
    and all time steps is either below a lower threshold or above an
    upper threshold:
    \begin{equation}
      \bar{a}_i < \alpha_{\text{lo}} \quad \text{or} \quad
      \bar{a}_i > \alpha_{\text{hi}}.
      \label{eq:method-dead-node}
    \end{equation}
  \item \textbf{Isolated:} it has neither incoming nor outgoing active
    connections.
  \item \textbf{Unreachable:} it is not on any path from an external
    input to an output node.
\end{enumerate}

Output nodes are never pruned. Furthermore, the network is protected
from shrinking below a minimum of \texttt{min\_hidden} hidden nodes; if
fewer hidden nodes would survive, the most active ones are retained.

Pruning is applied every $g_{\text{prune}}$ generations to the top
$\lfloor P/5 \rfloor$ networks. To avoid destroying functional networks,
pruning is only accepted if the resulting network has fitness at least
as high as the original (up to a small tolerance). This ``safe prune''
strategy ensures that pruning never degrades performance.

\begin{figure}[h]
  \centering
  \begin{tikzpicture}[
    every node/.style={font=\small},
    alive/.style={circle, draw, fill=blue!30, minimum size=0.6cm},
    dead/.style={circle, draw, fill=gray!30, minimum size=0.6cm,
      dashed},
    arrow/.style={-{Latex[length=1.5mm]}, thick},
  ]
    \node[alive] (a1) at (0, 2) {};
    \node[dead]  (a2) at (0, 1) {};
    \node[alive] (a3) at (0, 0) {};
    \node[alive] (b1) at (2, 2) {};
    \node[alive] (b2) at (2, 1) {};
    \node[dead]  (b3) at (2, 0) {};

    \draw[arrow] (a1) -- (a3);
    \draw[arrow] (a1) -- (b2);
    \draw[arrow] (a3) -- (b2);
    \draw[arrow] (b1) -- (a1);
    \draw[arrow] (b2) -- (a3);

    \node[font=\small] at (1, -1) {Before pruning};

    \draw[-{Latex[length=3mm]}, thick] (3.5, 1) -- (5, 1);

    \node[alive] (c1) at (6.5, 2) {};
    \node[alive] (c3) at (6.5, 0.5) {};
    \node[alive] (d1) at (8.5, 2) {};
    \node[alive] (d2) at (8.5, 1) {};

    \draw[arrow] (c1) -- (c3);
    \draw[arrow] (c1) -- (d2);
    \draw[arrow] (c3) -- (d2);
    \draw[arrow] (d1) -- (c1);
    \draw[arrow] (d2) -- (c3);

    \node[font=\small] at (7.5, -1) {After pruning};
  \end{tikzpicture}
  \caption{Pruning removes nodes whose activity is either below
    $\alpha_{\text{lo}}$ or above $\alpha_{\text{hi}}$ across all
    training samples (dashed gray nodes). The remaining nodes and
    connections are re-indexed to form a compact network.}
  \label{fig:method-prune}
\end{figure}

\subsection{Selection}
\label{sec:method-selection}

Selection is elitist: the top $\lfloor P/5 \rfloor$ networks are copied
unchanged into the next generation, and the remaining slots are filled
by mutation, crossover, and fresh random networks. The proportion of
fresh random networks is controlled by $p_{\text{fresh}}$, which we set
to $0.05$. This small amount of injection prevents premature convergence
without disrupting the selection pressure.

% ------------------------------------------------------------
\section{Complexity Analysis}
\label{sec:method-complexity}
% ------------------------------------------------------------

The computational cost of one evaluation of a network with $N$ nodes
and $T$ time steps is $O(N^2 T)$ in the worst case, because each node
may receive input from any other node. In practice, the connectivity is
sparse (Section~\ref{sec:method-prune}), so the cost is $O(ET)$, where
$E$ is the number of active connections. For the networks in our
experiments, $N \leq 27$ and $E \leq 357$, so a single evaluation takes
on the order of microseconds on a modern CPU.

The evolutionary algorithm performs $O(P G)$ evaluations, where $P$ is
the population size and $G$ is the number of generations. With $P = 300$
and $G = 1500$, this is $4.5 \times 10^5$ evaluations, which takes
approximately one minute on a single phone CPU.

% ------------------------------------------------------------
\section{Summary of Hyperparameters}
\label{sec:method-hyperparameters}
% ------------------------------------------------------------

Table~\ref{tab:method-hyperparameters} lists all hyperparameters used
in our experiments. We report the value used for most tasks; for tasks
with different values, the specific choice is noted in
Chapter~\ref{ch:experiments}.

\begin{table}[h]
  \centering
  \small
  \caption{Hyperparameters of \modnet{}.}
  \label{tab:method-hyperparameters}
  \begin{tabular}{lll}
    \toprule
    \textbf{Symbol} & \textbf{Description} & \textbf{Value} \\
    \midrule
    $N_0$ & Initial total nodes & 26 \\
    $T$ & Time steps per evaluation & 4--10 \\
    $P$ & Population size & 300 \\
    $G$ & Maximum generations & 400--1500 \\
    $p_{\text{cross}}$ & Crossover rate & 0.30 \\
    $p_{\text{fresh}}$ & Fresh random rate & 0.05 \\
    $\sigma_{\text{base}}$ & Base mutation scale & 0.30 \\
    $\sigma_{\max}$ & Maximum mutation scale & 1.50 \\
    $\gamma$ & Mutation growth factor & 1.25 \\
    $g_{\text{thresh}}$ & Plateau threshold & 25 \\
    $g_{\text{birth}}$ & Birth interval & 100 \\
    $g_{\text{prune}}$ & Prune interval & 20 \\
    $\epsilon_{\text{prune}}$ & Prune tolerance & $-0.001$ \\
    $\alpha_{\text{lo}}$ & Lower activity threshold & 0.02 \\
    $\alpha_{\text{hi}}$ & Upper activity threshold & 0.98 \\
    \texttt{min\_hidden} & Minimum hidden nodes & 2 \\
    $w_{\text{thresh}}$ & Weight pruning threshold & 0.04 \\
    \bottomrule
  \end{tabular}
\end{table}

% ------------------------------------------------------------
\section{Reproducibility}
\label{sec:method-reproducibility}
% ------------------------------------------------------------

The complete implementation of \modnet{}, including the evolutionary
algorithm, the fitness functions, the mutation, birth, and pruning
operators, and the experimental framework that generates all figures
and tables reported in this thesis, is available as open-source
software. The implementation depends only on NumPy and the Python
standard library; it requires no GPU and runs on commodity hardware,
including mobile phones.

All experiments are seeded, and the seed is reported in the experimental
setup (Chapter~\ref{ch:experiments}). The seed is reset per problem, so
that each benchmark runs with the same initial conditions regardless of
execution order. The output of each experiment is stored as JSON, CSV,
PNG, and plain-text files, which allows full reproduction and
inspection of every evolved network.

% ============================================================
%  Bibliography for this chapter
% ============================================================
